\chapter{Evaluating an expensive function in parallel}
\label{ch:parallel}

Every algorithm in the package needs the same service: evaluate a function at a batch of points,
as many at once as the cluster allows, and never lose or repeat an evaluation. This chapter
describes that shared layer. The arguments introduced here behave identically in
\code{slurm_minimize}, \code{slurm_mcmc} and \code{slurm_mcmc_hybrid}; \cref{app:common} collects
them in one table.

\section{The function you supply}

The expensive function takes a parameter vector --- a NumPy array of length $d$ --- and returns a
single number: a loss to be minimized (\cref{ch:optimization}), or a log-probability to be sampled
(\cref{ch:mcmc,ch:surrogate,ch:hybrid}), known up to a constant.

\begin{lstlisting}
def log_prob(theta):
    write_simulation_input(theta)      # e.g. a namelist for the forward code
    run_simulation()                   # minutes to hours
    return -0.5 * chi_squared(read_simulation_output())
\end{lstlisting}

Three facilities make real forward models easy to plug in.

\textbf{Its own working directory.} On the cluster each evaluation runs with its current directory
set to a directory of its own (\cref{sec:audit}), so the function can write input files, launch an
external code and parse its output without evaluations interfering with each other --- and every
file it writes stays behind as a record.

\textbf{A constant extra argument.} Anything the function needs besides $\theta$ --- the data, a
configuration dictionary, a path --- can be passed once as \code{extra_arg}. The function is then
called as \code{fun(theta, extra_arg)}.

\textbf{Deferred import.} A function has to be pickled to be sent to a compute node, and functions
defined in a script or notebook often cannot be. Instead of the function, pass where to find it:
\begin{lstlisting}
log_prob_fun = {'module_dir': '/home/me/model',
                'module_name': 'my_model',
                'function_name': 'log_prob'}
\end{lstlisting}
Only the three names are pickled. The function is imported inside each job, and re-imported fresh
there, so edits to the module are picked up by later evaluations. The same form is accepted for
\code{constraint_fun}.

\section{SlurmPool: one job array per batch}
\label{sec:slurmpool}

\code{SlurmPool} provides a single method, \code{map(fun, points)}, that returns the function's
value at each point, in order. It is a drop-in replacement for the \code{map} of
\code{multiprocessing.Pool} --- which is exactly how \code{emcee} uses it to evaluate its walkers.

Each call to \code{map} is submitted as \textbf{one Slurm job array}, named
\code{<job_name>_<call number>}, whose task $k$ evaluates point $k$. An iteration of an optimizer or
of MCMC is therefore a single scheduler transaction, however many points it contains. Batches larger
than \code{budget} points are split into consecutive arrays of at most that size.

\code{SlurmPool} can also be used on its own, for instance for a parameter scan:
\begin{lstlisting}
from slurmcmc.slurm_utils import SlurmPool

pool = SlurmPool(work_dir='scan', job_name='scan', cluster='slurm',
                 dim_input=3, dim_output=1,
                 submitit_kwargs={'slurm_partition': 'compute', 'timeout_min': 60})
values = pool.map(log_prob, points)
\end{lstlisting}

\section{Where the evaluations run}

The \code{cluster} argument selects the execution backend:

\begin{center}
\small
\begin{tabularx}{\linewidth}{@{}l L@{}}
  \toprule
  \code{cluster} & \textbf{behaviour} \\
  \midrule
  \code{'slurm'}     & submit job arrays to a Slurm cluster through submitit \\
  \code{'local'}     & run as local processes through submitit's local executor, with the same
                       directory layout as \code{'slurm'} --- useful for debugging the file handling \\
  \code{'local-map'} & evaluate sequentially inside the calling process: no jobs and no
                       directories. Fastest for cheap analytic functions and for tests \\
  \bottomrule
\end{tabularx}
\end{center}

A practical progression is to develop with \code{'local-map'}, check the files a real evaluation
writes with \code{'local'}, and then run with \code{'slurm'}. Scheduler options --- partition, time
limit, cores, memory --- are passed through \code{submitit_kwargs}, which is forwarded to submitit's
\code{AutoExecutor.update_parameters} (for example \code{slurm_partition}, \code{timeout_min},
\code{cpus_per_task}, \code{mem_gb}).

\section{The audit trail}
\label{sec:audit}

With \code{'slurm'} or \code{'local'}, every evaluation leaves a directory behind:

\begin{lstlisting}[style=plain]
work_dir/
  points_history.txt      every point of every call, led by call and point index
  values_history.txt     the corresponding results
  0/                  the first map() call, e.g. iteration 0
    inputs.txt        every point of the call, one per row
    outputs.txt       the corresponding results
    extra_arg.pkl     the extra_arg, when one was given
    0/
      input.txt       the point
      output.txt      its result
      ...             anything the function itself wrote
    1/
    ...
  1/                  the second call
  ...
\end{lstlisting}

The submitit logs of each job sit in the call's directory too, so a failed evaluation can be
debugged from its own inputs, outputs and error log, and a point's own output files can be re-used
later, for instance to compute a different quantity from the same runs. Once a call's results are
in, its points and results are also appended to \code{work_dir/points_history.txt} and
\code{work_dir/values_history.txt}, one row per point, led by the call and point indices.

When only the results matter, the call directories are clutter. \code{keep_run_dirs} decides what
stays: \code{'all'} (the default) keeps them, \code{'failed'} removes those in which every
evaluation succeeded, and \code{'none'} removes them all. The two history files are written in
every mode. A fresh run expects \code{work_dir} not to hold an earlier run's output, numbered call
directories or \code{points_history.txt}; use a new \code{work_dir}, or resume the old run from its
restart file (\cref{sec:restart}).

\section{Failures}
\label{sec:failures}

On a cluster some evaluations fail: a node dies, a simulation diverges, a job runs out of time.
An evaluation is counted as \textbf{failed} when
\begin{itemize}
  \item it returns \code{None}, NaN, or \code{job_fail_value};
  \item its job raises an exception; or
  \item it stays in the running state longer than \code{check_output_timeout_minutes}, in which case
        its job is cancelled.
\end{itemize}
Its result is then recorded as \code{job_fail_value}, and the run continues. The default is NaN for
minimization and $-10^{10}$ for the MCMC functions: a walker proposed into a failed point sees an
extremely low log-probability and simply rejects the move. Failed \emph{submissions} --- a
scheduler hiccup rather than a failed evaluation --- are retried up to
\code{submit_retry_max_attempts} times, \code{submit_retry_wait_seconds} apart;
\code{submit_delay_seconds} adds a pause after each submission, and
\code{check_output_interval_seconds} sets how often job states are polled.

\section{The evaluation history}

Every point a pool has evaluated is recorded. The pool, returned in each function's status as
\code{status['slurm_pool']}, exposes:

\begin{center}
\small
\begin{tabularx}{\linewidth}{@{}l L@{}}
  \toprule
  \textbf{attribute} & \textbf{content} \\
  \midrule
  \code{points_history}       & all evaluated points, in submission order, shape $(n, d)$ \\
  \code{values_history}       & their results, shape $(n, 1)$; failed points hold \code{job_fail_value} \\
  \code{inds_success_points}, \code{inds_failed_points} & row indices of successful and failed evaluations \\
  \code{evaluated_points_set} & the set of evaluated points, for instant duplicate checks \\
  \code{point_loc_dict}       & for each point, the (call, index) directory holding its files \\
  \code{num_evaluated_points} & the total count \\
  \bottomrule
\end{tabularx}
\end{center}

The duplicate check is what lets \code{slurm_minimize} guarantee it never pays twice for the same
point, and the history is what a surrogate is trained on. For very long MCMC runs whose history is
not needed, \code{record_history=False} saves the memory --- and time, since the bookkeeping grows
with the number of evaluations. It governs these in-memory arrays only: the same record is written
to disk as it goes, in \code{points_history.txt} and \code{values_history.txt} (\cref{sec:audit}).

\section{Restarting an interrupted run}
\label{sec:restart}

A calibration can run for days, and jobs get killed. With \code{save_restart=True} the complete
state of a run --- optimizer or sampler, evaluation history, iteration counter --- is written to
\code{work_dir/restart_file} after every \code{restart_save_interval} iterations. The write is
atomic: the file is written under a temporary name and then renamed, so a job killed mid-save
leaves the previous restart file intact.

A \code{SlurmPool} used on its own has no restart file of its own; it takes \code{load_restart=True}
instead, which rebuilds its call counter and its history from the two history files and continues
the numbering, evaluating whatever the next \code{map()} is given.

To continue, call the same function again with \code{load_restart=True}. It resumes from the saved
state and runs \code{num_iters} \emph{further} iterations, so the same call also extends a run that
finished. The default file names are \code{opt_restart.pkl}, \code{mcmc_restart.pkl} and
\code{hybrid_restart.pkl}; \code{slurm_mcmc} and \code{slurm_mcmc_hybrid} also accept the state as
an in-memory dictionary, \code{status_restart}.

By default each function installs a handler for \code{SIGTERM}, so that a run cancelled with
\code{scancel} exits with an error code and Slurm marks it as failed rather than completed. Pass
\code{install_signal_handler=False} when embedding a run in an application that manages its own
signals.

\section{Reproducible runs}
\label{sec:seed}

Every algorithm in the package makes random choices: the proposals of an MCMC sampler, the
candidates of an optimizer, the validation points and model fits of the hybrid pipeline. Pass
\code{random_seed} to make a run reproducible:
\begin{lstlisting}
status = slurm_mcmc(..., random_seed=0)
\end{lstlisting}
The seed is applied at the start of the loop, inside the process that runs it --- not by the
script that calls the function. That distinction matters in remote mode (\cref{sec:remote}): the
loop then runs in a fresh process on a compute node, which a \code{np.random.seed} call in the
submitting script never reaches, so without \code{random_seed} a remote run follows a different
trajectory from the same configuration. It seeds numpy's global generator, from which \texttt{emcee},
\texttt{nevergrad} and scikit-learn draw, Python's own generator, and torch's when \texttt{botorch}
is in use.

The restart file stores the state of those generators, so a run that was interrupted and resumed
reproduces the run that was never interrupted, iteration for iteration.

\begin{cautionbox}[Reproducible on the same software stack]
A seed fixes every random choice, not the arithmetic. Different versions of numpy, scipy or
scikit-learn, or a different linear-algebra library or CPU, change results in the last digits, and a
long MCMC or a sequence of model fits amplifies those differences. Expect identical results from the
same environment, and statistically equivalent ones elsewhere. A function that is itself random, such
as a stochastic simulation, has to seed itself.
\end{cautionbox}

\section{Running the whole loop as a job: remote mode}
\label{sec:remote}

The loop that drives an optimization or MCMC --- asking for points, waiting for their jobs,
updating the algorithm --- is itself a long-running process. Login nodes usually limit how long a
process may run. With \code{remote=True} the whole loop is instead submitted as a job of its own:
its configuration is pickled and run on \code{remote_cluster} (\code{'slurm'} or \code{'local'}),
under the name \code{main_{job_name}}, with a 30-day time limit unless
\code{remote_submitit_kwargs} says otherwise.

The call then returns a submitit \code{Job} immediately, instead of the status dictionary.
\code{job.result()} waits for the run to finish and returns the same dictionary a local run would.
submitit also stores it in \code{work_dir} as \code{<job_id>_0_result.pkl}, so it can be recovered
later from nothing but the directory and the job id:
\begin{lstlisting}
import submitit
status = submitit.SlurmJob(folder='mcmc', job_id='123456').result()
\end{lstlisting}

\begin{keypoint}[What survives a remote run]
The result file is written only when the job finishes. A failed job stores its traceback there
instead, and \code{.result()} raises it; a job killed before the end leaves no result file at all,
but its restart file (\cref{sec:restart}) is still there to resume from.
\end{keypoint}

\section{Configuration objects}

Each entry point is a thin wrapper that packs its keyword arguments into a configuration dataclass
and runs it. The dataclasses and runners can be used directly, for instance to build, inspect or
store a configuration before running it:
\begin{lstlisting}
from slurmcmc.optimization import MinimizeConfig, Minimizer
from slurmcmc.mcmc import MCMCConfig, MCMCRunner
from slurmcmc.hybrid import HybridMCMCConfig, HybridMCMCRunner

status = MCMCRunner(MCMCConfig(log_prob_fun=log_prob, init_points=init_points,
                               num_iters=1000)).run()
\end{lstlisting}
